\documentclass[a4paper,14pt]{scrartcl}
\usepackage{../main}

\title{Numeryczne wyznaczanie pochodnej funkcji oraz analiza błędu przybliżenia}
\author{Franciszek Szarek}
\nrprojektu{1} 
\date{\today}

\begin{document}
\maketitle

\tableofcontents
\newpage

\section*{Wstęp teoretyczny}
Z definicji analizy matematycznej, pochodna funkcji jest granicą ilorazu różnicowego:
\begin{equation*}
    \frac{du}{dx} = \lim_{x_2 \to x_1} \frac{u_2 - u_1}{x_2 - x_1}
\end{equation*}
Ponieważ systemy komputerowe nie są w stanie operować na pojęciu granicy dążącej do zera, w metodach numerycznych odstępuje się od przejścia granicznego na rzecz stałego, dyskretnego kroku siatki $h = x_2 - x_1$. Wymusza to zamianę dokładnego znaku równości na znak przybliżenia ($\approx$), co z kolei prowadzi do powstania błędu numerycznego, który jest przedmiotem analizy w niniejszym projekcie. Oszacowanie tego błędu opiera się na rozwinięciu funkcji w szereg Taylora.

\section{Wykres funkcji}
\begin{figure}[htbp]
    \centering
    \includegraphics[width=\textwidth]{wykres_sinh.png}
    \caption{Wykres funkcji $u(x) = \sinh(x)$ na przedziale $[-1, 1]$.}
    \label{fig:wykres_sinh}
\end{figure}

\newpage

\section{Dyskretyzacja przedziału}

W celu numerycznego wyznaczenia pochodnych, ciągłą dziedzinę funkcji
$u(x) = \sinh(x)$, określoną na przedziale $[-1, 1]$, poddano procesowi dyskretyzacji.
Do obliczeń przyjęto siatkę równomierną składającą się z $n = 8$ węzłów.
W pierwszej kolejności wyznaczono stały krok siatki $h$, określający
odległość pomiędzy sąsiadującymi węzłami. Wykorzystano w tym celu zależność:

\begin{equation*}
    h = \frac{x_n - x_1}{n - 1}
\end{equation*}

Podstawiając granice przedziału ($x_1 = -1$, $x_n = 1$) oraz wybraną liczbę węzłów ($n = 8$),
otrzymano:

\begin{equation*}
    h = \frac{1 - (-1)}{8 - 1} = \frac{2}{7} \approx 0.2857
\end{equation*}

Następnie wygenerowano zbiór węzłów $\{x_{\nu}\}$ dla $\nu = 1, 2, \dots, n$,
korzystając z jawnego wzoru na kolejne punkty siatki:

\begin{equation*}
    x_{\nu} = x_1 + (\nu - 1)h
\end{equation*}

Dla przyjętych założeń zbiór ten przyjmuje następujące dokładne wartości:
\begin{equation*}
    \begin{aligned}
        \{x_\nu\} = \{ & -1.000000, -0.714286, -0.428571, -0.142857, \\
                       & 0.142857, 0.428571, 0.714286, 1.000000 \}
    \end{aligned}
\end{equation*}

Każdemu z wyznaczonych węzłów przyporządkowano odpowiadającą mu dyskretną wartość funkcji $\{u_{\nu}\}$, określoną wzorem:

\begin{equation*}
    u_{\nu} = \sinh(x_{\nu})
\end{equation*}

Tak zdefiniowana siatka punktów $\{(x_{\nu}, u_{\nu})\}$ stanowi punkt wyjścia do przybliżania wartości pochodnej funkcji w tych węzłach za pomocą ilorazów różnicowych.
$$\begin{aligned}
        \{(x_\nu, u_\nu)\} = \{ & (-1.000000, -1.175201), (-0.714286, -0.776593), \\
                                & (-0.428571, -0.441812), (-0.142857, -0.143344), \\
                                & (0.142857, 0.143344), (0.428571, 0.441812),     \\
                                & (0.714286, 0.776593), (1.000000, 1.175201) \}
    \end{aligned}$$

\newpage
\begin{figure}[htbp]
    \centering
    \includegraphics[width=\textwidth]{wykres_sinh_wezly.png}
    \caption{Wykres funkcji $u(x) = \sinh(x)$ na przedziale $[-1, 1]$ z zaznaczonymi węzłami z dyskretyzacji.}
    \label{fig:wykres_sinh_wezly}
\end{figure}

\newpage

\section{Ilorazy różnicowe pierwszego rzędu}
Komputerowe wyznaczanie pochodnych funkcji ciągłych wymaga zastosowania metod przybliżonych. Podstawowym narzędziem do tego służącym są ilorazy różnicowe, wyprowadzane na podstawie rozwinięcia funkcji w szereg Taylora z pominięciem reszty $R_n$ (tzw. błąd obcięcia).


\subsection{Wyprowadzenie ilorazów różnicowych z szeregu Taylora}
Podstawą matematyczną do wyprowadzenia wzorów na numeryczne przybliżenie pochodnych (ilorazów różnicowych) jest rozwinięcie funkcji w szereg Taylora.
Szereg ten pozwala na wyznaczenie wartości funkcji w dowolnym punkcie $x$ na podstawie znajomości jej wartości oraz pochodnych w pewnym punkcie odniesienia $a$.
Ogólny wzór analityczny szeregu Taylora przyjmuje postać:


\begin{equation*}
    f(x) = f(a) + \frac{x-a}{1!}f^{(1)}(a) + \frac{(x-a)^2}{2!}f^{(2)}(a) + \dots + \frac{(x-a)^n}{n!}f^{(n)}(a) + R_n
\end{equation*}

gdzie $R_n$ oznacza resztę szeregu (błąd przybliżenia). \\

W celu zaadaptowania tego wzoru na potrzeby Metody Różnic Skończonych (FDM), wprowadzono następujące oznaczenia i podstawienia wynikające z przyjętej siatki węzłów:

\begin{enumerate}
    \item \textbf{Punkt odniesienia $a$} zastąpiono współrzędną węzła bieżącego: $x_\nu$
    \item \textbf{Wartość funkcji $f(a)$} w tym węźle oznaczono jako $u_\nu$
    \item \textbf{Punkt docelowy $x$} zastąpiono współrzędną węzła następnego: $x_{\nu+1}$
    \item \textbf{Wartość funkcji $f(x)$} w węźle następnym oznaczono jako $u_{\nu+1}$
    \item \textbf{Przyrost argumentu $(x - a)$}, będący odległością pomiędzy sąsiadującymi węzłami ($x_{\nu+1} - x_\nu$), zastąpiono stałym krokiem siatki: $h$.
    \item \textbf{Kolejne rzędy pochodnych $f^{(1)}(a), f^{(2)}(a), f^{(3)}(a)$} zapisano w uproszczonej notacji indeksowej jako $u'_\nu, u''_\nu, u'''_\nu$.
\end{enumerate}

\newpage

Podstawiając powyższe założenia do wzoru ogólnego, otrzymujemy równanie opisujące wartość funkcji w węźle następnym ($x_{\nu+1}$) w oparciu o parametry węzła bieżącego ($x_\nu$):

\begin{equation}
    \label{eq:szereg_taylora_progresywny}
    u_{\nu+1} = u_\nu + \frac{h}{1!}u'_\nu + \frac{h^2}{2!}u''_\nu + \frac{h^3}{3!}u'''_\nu + \dots + R_n
\end{equation}

Analogicznie, podstawiając za przyrost argumentu wartość ujemną ($-h$), uzyskujemy równanie dla węzła poprzedniego ($x_{\nu-1}$):

\begin{equation}
    u_{\nu-1} = u_\nu - \frac{h}{1!}u'_\nu + \frac{h^2}{2!}u''_\nu - \frac{h^3}{3!}u'''_\nu + \dots + R_n
\end{equation}

\subsubsection{Iloraz różnicowy progresywny (w przód)}
Do przybliżenia pochodnej w węźle $x_{\nu}$ wykorzystywana jest wartość funkcji w tym węźle oraz w węźle następnym ($x_{\nu+1}$).
Korzystając z wcześniej wyprowadzonego wzoru \eqref{eq:szereg_taylora_progresywny} aby wyznaczyć szukaną pierwszą pochodną ($u'_\nu$), w pierwszej kolejności przenosimy wartość funkcji z węzła bieżącego ($u_\nu$) na lewą stronę równania:
\begin{equation*}
    u_{\nu+1} - u_\nu = h \cdot u'_\nu + \frac{h^2}{2}u''_\nu + \dots + R_n
\end{equation*}


Następnie dzielimy obustronnie równanie przez krok siatki $h$, aby odizolować pochodną
$u'_\nu$:$$\frac{u_{\nu+1} - u_\nu}{h} = u'_\nu + \frac{h}{2}u''_\nu + \dots + \frac{R_n}{h}$$
Zamieniając stronami dla większej czytelności, otrzymujemy dokładną wartość pochodnej:
$$u'_\nu = \frac{u_{\nu+1} - u_\nu}{h} - \left( \frac{h}{2}u''_\nu + \dots + \frac{R_n}{h} \right)$$

W metodach numerycznych, aby uzyskać wzór obliczeniowy, odrzuca się wyrażenia zawierające wyższe rzędy pochodnych oraz resztę szeregu. Stanowią one tzw. błąd obcięcia. Po ich pominięciu równanie dokładne staje się równaniem przybliżonym:
$$u'_\nu \approx \frac{u_{\nu+1} - u_\nu}{h}$$

\newpage
\subsubsection{Iloraz różnicowy regresywny (wstecz)}
Analogicznie, przekształcamy drugie równanie szeregu Taylora (dla węzła poprzedniego), izolując $u'_{\nu}$:
\begin{equation*}
    h u'_{\nu} = u_{\nu} - u_{\nu-1} + \frac{h^2}{2!}u''_{\nu} - \dots
\end{equation*}

Po podzieleniu przez $h$ i odrzuceniu błędu obcięcia, uzyskujemy iloraz różnicowy regresywny:
\begin{equation*}
    u'_\nu \approx \frac{u_\nu - u_{\nu-1}}{h}
\end{equation*}

\newpage
\subsubsection{Iloraz różnicowy centralny}
Wyprowadzenie ilorazu centralnego opiera się na jednoczesnym wykorzystaniu rozwinięć w szereg Taylora zarówno dla węzła następnego ($x_{\nu+1}$), jak i poprzedniego ($x_{\nu-1}$). Zastosowanie tego podejścia pozwala na redukcję błędu obcięcia i uzyskanie wyższej dokładności numerycznej.

\subsubsection*{Wyprowadzenie wzoru}
Zapisujemy rozwinięcia w szereg Taylora dla kroku w przód i w tył, rozwijając je dla precyzji aż do trzeciej pochodnej:

\begin{equation}
    \label{eq:wzory_centralny_1}
    u_{\nu+1} = u_\nu + h \cdot u'_\nu + \frac{h^2}{2}u''_\nu + \frac{h^3}{6}u'''_\nu + \dots
\end{equation}

\begin{equation}
    \label{eq:wzory_centralny_2}
    u_{\nu-1} = u_\nu - h \cdot u'_\nu + \frac{h^2}{2}u''_\nu - \frac{h^3}{6}u'''_\nu + \dots
\end{equation}
Aby wyeliminować wartość funkcji w węźle bieżącym ($u_\nu$) oraz człony zawierające drugą pochodną ($u''_\nu$), odejmujemy równanie \eqref{eq:wzory_centralny_2} od równania \eqref{eq:wzory_centralny_1} stronami:

\begin{equation*}
    u_{\nu+1} - u_{\nu-1} = (u_\nu - u_\nu) + (h \cdot u'_\nu - (-h \cdot u'_\nu)) + \left(\frac{h^2}{2}u''_\nu - \frac{h^2}{2}u''_\nu\right) + \left(\frac{h^3}{6}u'''_\nu - \left(-\frac{h^3}{6}u'''_\nu\right)\right)
\end{equation*}


Po zredukowaniu wyrazów podobnych otrzymujemy: $$u_{\nu+1} - u_{\nu-1} = 2h \cdot u'_\nu + 2\frac{h^3}{6}u'''_\nu + \dots$$
co po uproszczeniu ułamka daje: $$u_{\nu+1} - u_{\nu-1} = 2h \cdot u'_\nu + \frac{h^3}{3}u'''_\nu + \dots$$

Następnie izolujemy szukaną pierwszą pochodną, dzieląc obustronnie równanie przez
$2h$:$$u'_\nu = \frac{u_{\nu+1} - u_{\nu-1}}{2h} - \frac{h^2}{6}u'''_\nu - \dots$$

Odrzucając końcowe wyrażenia szeregu stanowiące błąd obcięcia, uzyskujemy ostateczny wzór obliczeniowy:
$$u'_\nu \approx \frac{u_{\nu+1} - u_{\nu-1}}{2h}$$


\newpage
\clearpage
\vspace*{\fill}
\begin{center}
    \section{Zastosowanie ilorazów różnicowych do wyznaczenia pochodnej funkcji}
\end{center}
\vspace*{\fill}
\clearpage

\newpage

\subsection{Iloraz różnicowy progresywny}

\begin{table}[h!]
    \centering
    \small
    \renewcommand{\arraystretch}{1.2}
    \begin{tabular}{|c|c|c|c|}
        \hline
        \textbf{Węzeł $x$} & \textbf{Dokładna $u'$} & \textbf{Numeryczna $\tilde{u}'$} & \textbf{Błąd [\%]} \\ \hline
        $-1.00$            & $    1.5431$           & $    1.3951$                     & $ 9.59\%$          \\ \hline
        $-0.71$            & $    1.2661$           & $    1.1717$                     & $ 7.46\%$          \\ \hline
        $-0.43$            & $    1.0933$           & $    1.0446$                     & $ 4.45\%$          \\ \hline
        $-0.14$            & $    1.0102$           & $    1.0034$                     & $ 0.67\%$          \\ \hline
        $ 0.14$            & $    1.0102$           & $    1.0446$                     & $ 3.41\%$          \\ \hline
        $ 0.43$            & $    1.0933$           & $    1.1717$                     & $ 7.18\%$          \\ \hline
        $ 0.71$            & $    1.2661$           & $    1.3951$                     & $10.19\%$          \\ \hline
    \end{tabular}
\end{table}


\begin{figure}[htbp]
    \centering
    \includegraphics[width=\textwidth]{wykres_progresywna.png}
\end{figure}

Wyniki dla metody progresywnej (w przód) wykazują silną asymetrię i rosnącą niedokładność wzdłuż osi. Największy błąd ($10.19\%$) występuje w tym przypadku po prawej stronie dziedziny. Ponieważ metoda ta opiera się wyłącznie na danych "przyszłych" (węzłach następnych) i całkowicie ignoruje to, jak funkcja zachowywała się przed badanym punktem, generuje ona duże przekłamania wewnątrz osi. Z tego powodu w praktyce schemat ten stosuje się wyłącznie z konieczności na samym lewym brzegu siatki (na początku dziedziny), gdzie brakuje nam wcześniejszych punktów, aby móc użyć dokładniejszej metody centralnej.

\newpage
\subsection{Iloraz różnicowy regresywny}
\begin{table}[h!]
    \centering
    \small
    \renewcommand{\arraystretch}{1.2}
    \begin{tabular}{|c|c|c|c|}
        \hline
        \textbf{Węzeł $x$} & \textbf{Dokładna $u'$} & \textbf{Numeryczna $\tilde{u}'$} & \textbf{Błąd [\%]} \\ \hline
        $-0.71$            & $    1.2661$           & $    1.3951$                     & $10.19\%$          \\ \hline
        $-0.43$            & $    1.0933$           & $    1.1717$                     & $ 7.18\%$          \\ \hline
        $-0.14$            & $    1.0102$           & $    1.0446$                     & $ 3.41\%$          \\ \hline
        $ 0.14$            & $    1.0102$           & $    1.0034$                     & $ 0.67\%$          \\ \hline
        $ 0.43$            & $    1.0933$           & $    1.0446$                     & $ 4.45\%$          \\ \hline
        $ 0.71$            & $    1.2661$           & $    1.1717$                     & $ 7.46\%$          \\ \hline
        $ 1.00$            & $    1.5431$           & $    1.3951$                     & $ 9.59\%$          \\ \hline
    \end{tabular}

\end{table}

\begin{figure}[htbp]
    \centering
    \includegraphics[width=\textwidth]{wykres_regresywna.png}
\end{figure}
Wyniki dla metody regresywnej (w tył) stanowią niemal idealne, lustrzane odbicie błędów z poprzedniej tabeli. Jest to całkowicie naturalne zjawisko, ponieważ obie metody analizują dokładnie te same odcinki między punktami, ale "patrzą" na nie z przeciwnych kierunków. Największy błąd ($10.19\%$) występuje w tym przypadku po lewej stronie dziedziny. Ponieważ metoda ta opiera się na danych historycznych (węzłach poprzedzających) i całkowicie ignoruje to, jak funkcja zachowuje się dalej, generuje ona duże przekłamania wewnątrz osi. Z tego powodu w praktyce schemat ten stosuje się wyłącznie z konieczności na samym prawym brzegu siatki
(analogicznie jak w przypadku metody regresywnej, ale w przeciwnym kierunku).

\newpage
\subsection{Iloraz różnicowy centralny}
\begin{table}[h!]
    \centering
    \small
    \renewcommand{\arraystretch}{1.2}
    \begin{tabular}{|c|c|c|c|}
        \hline
        \textbf{Węzeł $x$} & \textbf{Dokładna $u'$} & \textbf{Numeryczna $\tilde{u}'$} & \textbf{Błąd [\%]} \\ \hline
        $-0.71$            & $    1.2661$           & $    1.2834$                     & $ 1.37\%$          \\ \hline
        $-0.43$            & $    1.0933$           & $    1.1082$                     & $ 1.37\%$          \\ \hline
        $-0.14$            & $    1.0102$           & $    1.0240$                     & $ 1.37\%$          \\ \hline
        $ 0.14$            & $    1.0102$           & $    1.0240$                     & $ 1.37\%$          \\ \hline
        $ 0.43$            & $    1.0933$           & $    1.1082$                     & $ 1.37\%$          \\ \hline
        $ 0.71$            & $    1.2661$           & $    1.2834$                     & $ 1.37\%$          \\ \hline
    \end{tabular}
\end{table}

\begin{figure}[htbp]
    \centering
    \includegraphics[width=\textwidth]{wykres_centralna.png}
\end{figure}

Wyniki z tabeli wyraźnie pokazują, że metoda centralna jest znacznie dokładniejsza niż schematy jednostronne. Błąd obliczeń spada tutaj do poziomu zaledwie $1.37\%$. Tak wysoka precyzja wynika z faktu, że metoda ta wyznaczając pochodną uwzględnia wartości funkcji po obu stronach badanego węzła (zarówno z przodu, jak i z tyłu). Dzięki temu symetrycznemu podejściu, większość niedokładności matematycznych wzajemnie się znosi, co drastycznie poprawia wynik końcowy.
Metody centralnej nie stosuje się na brzegach dziedziny, ponieważ do wykonania obliczeń wymaga ona punktów sąsiadujących z obu stron, których na skrajnych krawędziach brakuje.
\newpage

\section{ Analiza błędu względnego}

W celu ilościowej oceny dokładności zastosowanych metod numerycznych, dla każdego węzła siatki obliczono błąd względny $E$ wyrażony w procentach. Błąd ten definiuje się jako stosunek wartości bezwzględnej różnicy między pochodną dokładną (analityczną) $u'_{\nu}$, a jej przybliżeniem numerycznym $\tilde{u}'_{\nu}$, do wartości pochodnej dokładnej:

\begin{equation*}
    E = \left| \frac{u'_{\nu} - \tilde{u}'_{\nu}}{u'_{\nu}} \right| \cdot 100\%
\end{equation*}

W analizowanym przypadku funkcją bazową jest $u(x) = \sinh(x)$, zatem pochodna dokładna w powyższym wzorze przyjmuje postać $u'(x) = \cosh(x)$.

Obliczenia błędu względnego przeprowadzono dla wszystkich trzech zdefiniowanych wcześniej ilorazów różnicowych. Zestawienie wyników na wspólnym wykresie pozwala zaobserwować zachowanie każdej z metod. Metoda oparta na ilorazie centralnym generuje najmniejszy błąd przybliżenia, co wynika z faktu, iż w rozwinięciu w szereg Taylora redukcji ulegają składniki rzędu $h$, przez co metoda ta jest wyższego rzędu dokładności. Z kolei błędy dla ilorazu progresywnego i regresywnego są zauważalnie większe i rosną symetrycznie w kierunku przeciwległych brzegów przedziału.

\begin{figure}[htbp]
    \centering
    \includegraphics[width=0.8\textwidth]{wykres_bledy.png}
    \caption{Porównanie błędu względnego $E$ [\%] dla ilorazów różnicowych: progresywnego, regresywnego i centralnego.}
    \label{fig:wykres_bledy}
\end{figure}
\end{document}